\chapter{Hamiltonkreise}

\section{Dieselbe Frage, anderes Objekt}

Kapitel 6 fragte: Kann man jede \kursiv{Kante} genau einmal durchlaufen? Jetzt fragen wir: Kann man jede \kursiv{Ecke} genau einmal besuchen?

Die beiden Fragen sehen einander so ähnlich, dass man erwartet, die Antworten würden einander ähneln. Sie tun es nicht. Der Unterschied ist einer der lehrreichsten der ganzen Mathematik.

\begin{definition}
\boxtitelm{Definition 7.1 --- Hamiltonweg, Hamiltonkreis}
Sei $G$ ein Graph.
\begin{itemize}[leftmargin=1.4em, itemsep=4pt]
\item Ein \Begriff{Hamiltonweg} ist ein Weg, der jede Ecke von $G$ enthält.
\item Ein \Begriff{Hamiltonkreis} ist ein Kreis, der jede Ecke von $G$ enthält.
\item $G$ heißt \Begriff{hamiltonsch}, wenn er einen Hamiltonkreis besitzt.
\end{itemize}
\end{definition}

\begin{merksatz}
\boxtitel{Ein erster Vergleich}
\begin{center}
\begin{tikzpicture}[scale=1.15]
\begin{scope}
\foreach \i in {0,...,4} \node[knoten] (a\i) at (90+72*\i:1.0) {};
\foreach \i/\j in {0/1,1/2,2/3,3/4,4/0} \draw[kanted] (a\i) -- (a\j);
\foreach \i/\j in {0/2,1/3,2/4,3/0,4/1} \draw[kante] (a\i) -- (a\j);
\node[font=\scriptsize\itshape, color=grau] at (0,-1.6) {$K_5$: hamiltonsch \kursiv{und} eulersch};
\end{scope}
\begin{scope}[xshift=5.2cm]
\node[knoten] (z) at (0,0) {};
\node[knoten] (p1) at (-1.0,0.7) {};
\node[knoten] (p2) at (1.0,0.7) {};
\node[knoten] (p3) at (0,-1.1) {};
\draw[kante] (z)--(p1); \draw[kante] (z)--(p2); \draw[kante] (z)--(p3);
\node[font=\scriptsize\itshape, color=grau] at (0,-1.6) {Stern: weder noch};
\end{scope}
\end{tikzpicture}
\end{center}

\begin{center}
\begin{tikzpicture}[scale=1.15]
\begin{scope}
\node[knoten] (c1) at (-1.0,0.6) {};
\node[knoten] (c2) at (0,0.6) {};
\node[knoten] (c3) at (1.0,0.6) {};
\node[knoten] (c4) at (0,-0.6) {};
\draw[kante] (c1)--(c2); \draw[kante] (c2)--(c3);
\draw[kante] (c1)--(c4); \draw[kante] (c3)--(c4);
\draw[kante] (c2)--(c4);
\node[font=\scriptsize\itshape, color=grau] at (0,-1.3) {hamiltonsch, nicht eulersch};
\end{scope}
\begin{scope}[xshift=5.4cm]
\node[knoten] (m) at (0,0) {};
\node[knoten] (u1) at (-1.0,0.75) {};
\node[knoten] (u2) at (-1.0,-0.75) {};
\node[knoten] (v1) at (1.0,0.75) {};
\node[knoten] (v2) at (1.0,-0.75) {};
\draw[kante] (m)--(u1); \draw[kante] (m)--(u2); \draw[kante] (u1)--(u2);
\draw[kante] (m)--(v1); \draw[kante] (m)--(v2); \draw[kante] (v1)--(v2);
\node[font=\scriptsize\itshape, color=grau] at (0,-1.3) {eulersch, nicht hamiltonsch};
\end{scope}
\end{tikzpicture}
\end{center}

Alle vier Kombinationen kommen vor. Die beiden Begriffe sind also voneinander unabhängig --- keiner folgt aus dem anderen.
\end{merksatz}

\begin{warnung}
\boxtitelr{Der entscheidende Unterschied --- und warum er so groß ist}
\begin{center}
{\small
\renewcommand{\arraystretch}{1.5}
\begin{tabular}{@{}P{2.4cm}P{2.9cm}P{2.9cm}@{}}
\toprule
& \textbf{Eulerkreis} & \textbf{Hamiltonkreis}\\
\midrule
Kriterium & \textcolor{gruen}{Satz 6.4, vollständig} & \textcolor{akzent}{\textbf{keines bekannt}}\\
Prüfaufwand & \textcolor{gruen}{Grade zählen} & \textcolor{akzent}{im Wesentlichen Probieren}\\
Bei $n = 100$ & \textcolor{gruen}{Sekundenbruchteil} & \textcolor{akzent}{unlösbar im Allgemeinen}\\
\bottomrule
\end{tabular}}
\end{center}

Woher der Unterschied? Eine Antwort gibt der Beweis von Satz 6.4. Er funktioniert, weil sich Teilkreise \kursiv{verschmelzen} lassen: Zwei kantendisjunkte geschlossene Züge mit gemeinsamer Ecke ergeben zusammen wieder einen. Man kann also lokal arbeiten und die Stücke zusammensetzen.

Bei Hamiltonkreisen geht das nicht. Zwei Kreise mit gemeinsamer Ecke ergeben keinen Kreis, sondern eine Acht. Es gibt kein Verschmelzungsargument, und deshalb hilft lokales Arbeiten nicht.

\begin{center}
\begin{tikzpicture}[scale=1.05]
\begin{scope}
\draw[kante] (-0.75,0) circle (0.75);
\draw[kanted] (0.75,0) circle (0.75);
\node[knotenr] (x) at (0,0) {};
\node[font=\scriptsize\itshape, color=gruen] at (0,-1.35) {Kantenzüge: verschmelzen zu einem};
\end{scope}
\begin{scope}[xshift=5.6cm]
\draw[kante] (-0.75,0) circle (0.75);
\draw[kanted] (0.75,0) circle (0.75);
\node[knotenr] (y) at (0,0) {};
\node[font=\scriptsize\itshape, color=akzent] at (0,-1.35) {Kreise: die Mitte hat Grad $4$};
\end{scope}
\end{tikzpicture}
\end{center}

\begin{center}
\begin{minipage}{0.88\textwidth}
\kursiv{Dass kein einfaches Kriterium bekannt ist, ist nicht bloß ein Wissenslücke. Kapitel~15 erklärt, warum es aus tiefen Gründen unwahrscheinlich ist, dass eines existiert.}
\end{minipage}
\end{center}
\end{warnung}

\section{Ein Werkzeug zum Widerlegen}

Wenn es kein Kriterium gibt, wie zeigt man dann, dass ein Graph \kursiv{nicht} hamiltonsch ist? Der folgende Satz liefert die meistgenutzte Antwort.

\begin{merksatz}
\boxtitel{Satz 7.2 --- die Komponentenbedingung}
Sei $G$ hamiltonsch. Dann gilt für jede nichtleere Teilmenge $S \subseteq V$:
\[ c(G - S) \;\le\; |S| , \]
wobei $c(\cdot)$ die Anzahl der Zusammenhangskomponenten bezeichnet.
\end{merksatz}

\begin{merksatz}
\boxtitel{Beweis}
\schritt{Schritt 1 --- Den Hamiltonkreis zerschneiden}
Sei $H$ ein Hamiltonkreis von $G$ und $S \subseteq V$ nichtleer mit $|S| = k$.

Wir entfernen die $k$ Ecken von $S$ aus $H$. Ein Kreis, aus dem man $k$ Ecken entfernt, zerfällt in höchstens $k$ Wegstücke.

\schritt{Warum höchstens $k$}
Denken Sie an einen Ring aus Perlen. Nimmt man eine Perle heraus, wird der Ring zu einer Kette --- ein Stück. Nimmt man eine zweite heraus, zerfällt die Kette in zwei Stücke. Jede weitere entfernte Perle erhöht die Stückzahl um höchstens eins.

Genauer: Nach dem Entfernen der ersten Ecke haben wir $1$ Stück. Jede weitere entfernte Ecke liegt in genau einem Stück und teilt dieses in höchstens $2$ Teile --- weniger, wenn sie am Rand liegt. Also nach $k$ Entfernungen höchstens $k$ Stücke.

\begin{center}
\begin{tikzpicture}[scale=1.15]
\foreach \i in {0,...,7} {
  \pgfmathsetmacro{\ang}{90-45*\i}
  \ifnum\i=1 \node[knotenr] (h\i) at (\ang:1.15) {}; \else
  \ifnum\i=5 \node[knotenr] (h\i) at (\ang:1.15) {}; \else
  \node[knoten] (h\i) at (\ang:1.15) {}; \fi \fi
}
\foreach \i/\j in {0/1,1/2,2/3,3/4,4/5,5/6,6/7,7/0} \draw[kante] (h\i) -- (h\j);
\node[font=\scriptsize\itshape, color=akzent] at (0,-1.75) {zwei Ecken heraus $\Rightarrow$ zwei Wegstücke};
\end{tikzpicture}
\end{center}

\schritt{Schritt 2 --- Von $H$ auf $G$ schließen}
$H$ ist ein Teilgraph von $G$ mit derselben Eckenmenge. Also ist $H - S$ ein Teilgraph von $G - S$ mit derselben Eckenmenge.

Fügt man zu einem Graphen Kanten hinzu, so kann die Komponentenzahl nur sinken oder gleich bleiben --- nie steigen. Denn zwei Ecken, die vorher verbunden waren, bleiben es.

Also
\[ c(G-S) \;\le\; c(H-S) \;\le\; k = |S| . \qquad \blacksquare \]
\end{merksatz}

\begin{merksatz}
\boxtitel{Wie man den Satz benutzt}
Der Satz sagt: \kursiv{Ist} $G$ hamiltonsch, \kursiv{dann} gilt die Ungleichung. Man benutzt ihn in der umgekehrten Richtung: Findet man ein $S$, für das sie verletzt ist, so ist $G$ nicht hamiltonsch.

\begin{center}
\begin{tikzpicture}[scale=1.2]
\node[knotenr] (m1) at (-0.6,0) {};
\node[knotenr] (m2) at (0.6,0) {};
\foreach \i/\y in {0/1.1, 1/0.0, 2/-1.1} {
  \node[knoten] (l\i) at (-1.9,\y) {};
  \node[knoten] (r\i) at (1.9,\y) {};
}
\node[knoten] (l3) at (-2.9,0.55) {};
\node[knoten] (r3) at (2.9,0.55) {};
\draw[kante] (m1)--(l0);\draw[kante] (m1)--(l1);\draw[kante] (m1)--(l2);
\draw[kante] (m2)--(r0);\draw[kante] (m2)--(r1);\draw[kante] (m2)--(r2);
\draw[kante] (m1)--(m2);
\draw[kante] (l0)--(l3);\draw[kante] (r0)--(r3);
\node[font=\scriptsize, color=akzent] at (0,-0.45) {$S$};
\end{tikzpicture}
\end{center}

\schritt{Beispiel}
Wählt man $S$ als die beiden roten Ecken, so zerfällt der Rest in sechs Teile: links drei (die drei Zweige, von denen einer zwei Ecken hat), rechts drei. Also
\[ c(G-S) = 6 \;>\; 2 = |S| . \]
Der Graph ist nicht hamiltonsch.

\begin{center}
\begin{minipage}{0.88\textwidth}
\kursiv{Als Sonderfall enthält der Satz: Ein hamiltonscher Graph hat keine Schnittecke. Denn für $S = \{v\}$ mit einer Schnittecke $v$ wäre $c(G-v) \ge 2 > 1$. Anschaulich: Ein Rundgang, der zurückkehren muss, kann keinen Engpass nur einmal durchqueren.}
\end{minipage}
\end{center}
\end{merksatz}

\begin{warnung}
\boxtitelr{Die Bedingung ist notwendig, aber nicht hinreichend}
Es gibt Graphen, die die Komponentenbedingung für \kursiv{jedes} $S$ erfüllen und trotzdem keinen Hamiltonkreis haben. Der Petersen-Graph ist so einer --- Abschnitt 7.5.

Das ist kein Mangel dieses Satzes, sondern die Regel: Alle bekannten notwendigen Bedingungen sind unvollständig, und alle bekannten hinreichenden Bedingungen sind zu stark. Dazwischen bleibt eine Lücke, die niemand schließen konnte.
\end{warnung}

\section{Zwei hinreichende Bedingungen}

Wenn schon kein vollständiges Kriterium, dann wenigstens Bedingungen, die den Hamiltonkreis \kursiv{garantieren}. Die berühmteste stammt von Ore; die bekanntere von Dirac folgt daraus.

\begin{merksatz}
\boxtitel{Satz 7.3 (Ore 1960)}
Sei $G$ ein Graph mit $n \ge 3$ Ecken. Gilt für je zwei \kursiv{nicht benachbarte} Ecken $u,v$
\[ \deg(u) + \deg(v) \;\ge\; n , \]
so ist $G$ hamiltonsch.
\end{merksatz}

\begin{merksatz}
\boxtitel{Beweis --- Schritt 1: den schlimmsten Fall betrachten}
Angenommen, der Satz wäre falsch. Dann gibt es einen Graphen $G$ auf $n \ge 3$ Ecken, der die Bedingung erfüllt und keinen Hamiltonkreis hat.

Wir fügen zu $G$ so lange Kanten hinzu, wie dabei kein Hamiltonkreis entsteht, und nennen das Ergebnis $G'$. Das Verfahren endet, weil man höchstens $\binom{n}{2}$ Kanten hat.

\schritt{Was $G'$ leistet}
\begin{center}
\renewcommand{\arraystretch}{1.5}
\begin{tabular}{@{}P{1.4cm}P{7.0cm}@{}}
\toprule
(a) & $G'$ erfüllt die Ore-Bedingung. Denn Hinzufügen von Kanten kann Grade nur erhöhen, und die Bedingung ist eine Mindestforderung.\\
(b) & $G'$ hat keinen Hamiltonkreis --- nach Konstruktion.\\
(c) & $G'$ ist nicht vollständig. Denn $K_n$ mit $n \ge 3$ hat einen Hamiltonkreis: Man läuft die Ecken in beliebiger Reihenfolge ab, alle nötigen Kanten sind vorhanden.\\
(d) & Fügt man zu $G'$ irgendeine fehlende Kante hinzu, so \kursiv{entsteht} ein Hamiltonkreis --- sonst hätte man weiter hinzugefügt.\\
\bottomrule
\end{tabular}
\end{center}
\end{merksatz}

\begin{merksatz}
\boxtitel{Beweis --- Schritt 2: ein Hamiltonweg entsteht}
Nach (c) gibt es zwei nicht benachbarte Ecken $u$ und $v$ in $G'$. Nach (d) hat $G' + \{u,v\}$ einen Hamiltonkreis $H$.

\schritt{$H$ muss die neue Kante benutzen}
Andernfalls läge $H$ ganz in $G'$, im Widerspruch zu (b).

\schritt{Die neue Kante herausnehmen}
$H$ ist ein Kreis durch alle $n$ Ecken. Entfernt man daraus die Kante $\{u,v\}$, so bleibt ein Weg durch alle Ecken, der bei $u$ beginnt und bei $v$ endet. Wir schreiben ihn
\[ u = v_1,\ v_2,\ v_3,\ \ldots,\ v_{n-1},\ v_n = v . \]
Das ist ein Hamiltonweg von $u$ nach $v$ in $G'$.
\end{merksatz}

\begin{merksatz}
\boxtitel{Beweis --- Schritt 3: der Zählkniff}
Jetzt kommt der Kern. Wir definieren zwei Indexmengen:
\[ S = \{\, i \;:\; 1 \le i \le n-1,\ \{u, v_{i+1}\} \in E(G') \,\}, \]
\[ T = \{\, i \;:\; 1 \le i \le n-1,\ \{v, v_{i}\} \in E(G') \,\} . \]

In Worten: $S$ sammelt die Stellen, an denen der \kursiv{Nachfolger} auf dem Weg ein Nachbar von $u$ ist; $T$ die Stellen, an denen die Ecke \kursiv{selbst} ein Nachbar von $v$ ist.

\schritt{Die Größen von $S$ und $T$}
Jeder Nachbar von $u$ ist eine der Ecken $v_2, \ldots, v_n$ (nicht $v_1 = u$ selbst). Schreibt man ihn als $v_{i+1}$, so ist $1 \le i \le n-1$. Verschiedene Nachbarn liefern verschiedene $i$. Also
\[ |S| = \deg(u) . \]
Ebenso ist jeder Nachbar von $v$ eine der Ecken $v_1, \ldots, v_{n-1}$, also
\[ |T| = \deg(v) . \]

\schritt{Beide liegen in derselben kleinen Menge}
$S$ und $T$ sind Teilmengen von $\{1, 2, \ldots, n-1\}$, einer Menge mit $n-1$ Elementen.

\schritt{Sie müssen sich überschneiden}
Wären $S$ und $T$ disjunkt, so wäre
\[ |S| + |T| = |S \cup T| \le n-1 . \]
Nach Voraussetzung ist aber $u \not\sim v$, also
\[ |S| + |T| = \deg(u) + \deg(v) \;\ge\; n \;>\; n - 1 . \]
Widerspruch. Also gibt es ein $i \in S \cap T$.
\end{merksatz}

\begin{merksatz}
\boxtitel{Beweis --- Schritt 4: der Kreis wird geschlossen}
Für dieses $i$ gilt:
\[ \{u, v_{i+1}\} \in E(G') \quad \text{und} \quad \{v, v_i\} \in E(G') . \]

Nun bauen wir einen Hamiltonkreis in $G'$:
\[ \underbrace{v_1 \to v_2 \to \cdots \to v_i}_{\text{vorwärts}} \;\to\; \underbrace{v_n \to v_{n-1} \to \cdots \to v_{i+1}}_{\text{rückwärts}} \;\to\; v_1 . \]

\begin{center}
\begin{tikzpicture}[scale=1.0]
\foreach \i/\x in {1/0, 2/1.1, 3/2.2, 4/3.3, 5/4.4, 6/5.5}
  \node[knoten] (w\i) at (\x,0) {};
\foreach \i/\j in {1/2,2/3,4/5,5/6} \draw[kante] (w\i) -- (w\j);
\draw[kante] (w3) -- (w4);
\node[font=\scriptsize, color=grau] at (0,-0.35) {$v_1=u$};
\node[font=\scriptsize, color=tief] at (2.2,-0.35) {$v_i$};
\node[font=\scriptsize, color=tief] at (3.3,-0.35) {$v_{i+1}$};
\node[font=\scriptsize, color=grau] at (5.5,-0.35) {$v_n=v$};
\draw[kanted] (w1) to[bend left=32] (w4);
\draw[kanted] (w3) to[bend right=32] (w6);
\node[font=\scriptsize, color=akzent] at (1.7,1.0) {$\{u,v_{i+1}\}$};
\node[font=\scriptsize, color=akzent] at (3.9,-1.0) {$\{v_i,v\}$};
\end{tikzpicture}
\end{center}

\schritt{Warum das ein Kreis durch alle Ecken ist}
Der erste Abschnitt durchläuft $v_1, \ldots, v_i$; der zweite $v_n, v_{n-1}, \ldots, v_{i+1}$. Zusammen sind das alle $n$ Ecken, jede genau einmal.

Die benutzten Kanten: die Wegkanten $\{v_1,v_2\}, \ldots, \{v_{i-1},v_i\}$ und $\{v_{i+1},v_{i+2}\}, \ldots, \{v_{n-1},v_n\}$ --- sie gehören zum Hamiltonweg. Dazu die Übergangskante $\{v_i, v_n\} = \{v_i, v\}$ und die Schlusskante $\{v_{i+1}, v_1\} = \{v_{i+1}, u\}$ --- beide liegen nach Wahl von $i$ in $G'$.

\schritt{Der Widerspruch}
$G'$ hat also einen Hamiltonkreis --- im Widerspruch zu (b). Die Annahme war falsch. $\blacksquare$
\end{merksatz}

\begin{merksatz}
\boxtitel{Folgerung 7.4 (Dirac 1952)}
Hat jede Ecke eines Graphen mit $n \ge 3$ Ecken mindestens den Grad $n/2$, so ist $G$ hamiltonsch.

\schritt{Beweis}
Für nicht benachbarte $u, v$ gilt
\[ \deg(u) + \deg(v) \;\ge\; \frac{n}{2} + \frac{n}{2} = n . \]
Das ist die Ore-Bedingung. $\blacksquare$

\begin{center}
\begin{minipage}{0.88\textwidth}
\kursiv{Historisch war es umgekehrt: Dirac kam acht Jahre früher, Ore verallgemeinerte. Dass die stärkere Aussage den kürzeren Beweis hat, ist in der Mathematik nicht selten --- eine schärfere Formulierung zeigt oft erst, worauf es wirklich ankommt.}
\end{minipage}
\end{center}
\end{merksatz}

\begin{warnung}
\boxtitelr{Die Schranke $n/2$ ist scharf --- und die Bedingung ist weit von notwendig entfernt}
\schritt{Scharf}
Nimmt man zwei vollständige Graphen $K_{(n+1)/2}$, die genau eine Ecke gemeinsam haben (für ungerades $n$), so hat jede Ecke mindestens Grad $\frac{n-1}{2}$ --- knapp unter $n/2$. Der Graph hat eine Schnittecke und ist nach Satz 7.2 nicht hamiltonsch.

\schritt{Weit von notwendig}
Der Kreis $C_{100}$ ist hamiltonsch --- er \kursiv{ist} sein eigener Hamiltonkreis. Jede Ecke hat aber nur Grad $2$, und $n/2 = 50$. Dirac ist also meilenweit verfehlt.

\begin{center}
\begin{minipage}{0.88\textwidth}
\kursiv{Das ist die Lage bei allen bekannten hinreichenden Bedingungen: Sie verlangen sehr dichte Graphen. Die schwierigen Fälle --- dünne Graphen wie der Petersen-Graph --- werden von keiner erfasst.}
\end{minipage}
\end{center}
\end{warnung}

\section{Der Würfelgraph und die Gray-Codes}

\begin{definition}
\boxtitelm{Definition 7.5 --- Würfelgraph}
Der \Begriff{Würfelgraph} $Q_d$ hat als Ecken alle $2^d$ Folgen aus Nullen und Einsen der Länge $d$. Zwei Ecken sind benachbart, wenn sie sich an genau einer Stelle unterscheiden.
\end{definition}

\begin{merksatz}
\boxtitel{Satz 7.6 --- $Q_d$ ist hamiltonsch für $d \ge 2$}
\schritt{Beweis durch Induktion über $d$}
\textbf{Anfang $d = 2$:} $Q_2$ hat die vier Ecken $00, 01, 11, 10$, und je zwei aufeinanderfolgende unterscheiden sich an einer Stelle. Das ist ein Viereck --- der Hamiltonkreis ist der Graph selbst.

\textbf{Schritt.} Sei $d \ge 3$ und $C$ ein Hamiltonkreis in $Q_{d-1}$.

\schritt{Die Zerlegung}
$Q_d$ besteht aus zwei Kopien von $Q_{d-1}$: der Menge aller Folgen, die mit $0$ beginnen, und der Menge aller, die mit $1$ beginnen. Innerhalb jeder Kopie sind die Nachbarschaften dieselben wie in $Q_{d-1}$ --- die erste Stelle ändert sich ja nicht. Zusätzlich ist jede Ecke $0w$ mit $1w$ benachbart, denn sie unterscheiden sich genau an der ersten Stelle.

\schritt{Der Zusammenbau}
Sei $\{x, y\}$ eine beliebige Kante des Hamiltonkreises $C$. Wir bilden:
\begin{center}
\renewcommand{\arraystretch}{1.45}
\begin{tabular}{@{}P{3.2cm}P{5.4cm}@{}}
\toprule
Kopie $0$ & der Kreis $C$, mit vorangestellter $0$, ohne die Kante $\{0x, 0y\}$ --- ein Weg von $0x$ nach $0y$\\
Kopie $1$ & derselbe Kreis mit vorangestellter $1$, ohne $\{1x, 1y\}$ --- ein Weg von $1x$ nach $1y$\\
Verbindung & die Kanten $\{0x, 1x\}$ und $\{0y, 1y\}$\\
\bottomrule
\end{tabular}
\end{center}

\begin{center}
\begin{tikzpicture}[scale=1.15]
\draw[kante] (-1.3,0) ellipse (0.95 and 1.15);
\draw[kante] (1.3,0) ellipse (0.95 and 1.15);
\node[knotenb] (a) at (-0.55,0.5) {};
\node[knotenb] (b) at (-0.55,-0.5) {};
\node[knotenb] (c) at (0.55,0.5) {};
\node[knotenb] (d) at (0.55,-0.5) {};
\draw[kanted] (a)--(c); \draw[kanted] (b)--(d);
\node[font=\scriptsize, color=grau] at (-0.95,0.78) {$0x$};
\node[font=\scriptsize, color=grau] at (-0.95,-0.78) {$0y$};
\node[font=\scriptsize, color=grau] at (0.95,0.78) {$1x$};
\node[font=\scriptsize, color=grau] at (0.95,-0.78) {$1y$};
\node[font=\scriptsize\itshape, color=grau] at (-1.3,-1.5) {Kopie $0$};
\node[font=\scriptsize\itshape, color=grau] at (1.3,-1.5) {Kopie $1$};
\end{tikzpicture}
\end{center}

\schritt{Warum das ein Hamiltonkreis ist}
Wir laufen: von $0x$ über den Weg in Kopie $0$ nach $0y$, dann über $\{0y, 1y\}$ nach $1y$, dann über den Weg in Kopie $1$ nach $1x$, dann über $\{1x, 0x\}$ zurück.

Jede Ecke von Kopie $0$ kommt genau einmal vor (der Weg ist ein Hamiltonweg dieser Kopie), ebenso jede von Kopie $1$. Zusammen alle $2^d$ Ecken, jede genau einmal. $\blacksquare$
\end{merksatz}

\begin{merksatz}
\boxtitel{Was das praktisch bedeutet: Gray-Codes}
Ein Hamiltonkreis in $Q_d$ ist dasselbe wie eine Anordnung aller $2^d$ Bitfolgen der Länge $d$, bei der sich aufeinanderfolgende Folgen an genau einer Stelle unterscheiden --- und die letzte wieder an einer Stelle von der ersten. So etwas heißt \Begriff{Gray-Code}.

\schritt{Der Fall $d = 3$, aus dem Beweis gewonnen}
$Q_2$-Kreis: $00 \to 01 \to 11 \to 10 \to 00$. Wir entfernen die Kante $\{10, 00\}$ und erhalten den Weg $00 \to 01 \to 11 \to 10$.

\begin{center}
\renewcommand{\arraystretch}{1.4}
\begin{tabular}{@{}P{4.2cm}P{4.2cm}@{}}
\toprule
\textbf{Kopie $0$ (vorwärts)} & \textbf{Kopie $1$ (rückwärts)}\\
\midrule
$000, 001, 011, 010$ & $110, 111, 101, 100$\\
\bottomrule
\end{tabular}
\end{center}

Der Gesamtkreis lautet
\[ 000 \to 001 \to 011 \to 010 \to 110 \to 111 \to 101 \to 100 \to 000 . \]

\schritt{Probe}
Wir prüfen jeden Übergang auf genau eine geänderte Stelle:
\begin{center}
{\small
\renewcommand{\arraystretch}{1.32}
\begin{tabular}{@{}P{2.4cm}P{1.8cm}P{2.4cm}P{1.8cm}@{}}
\toprule
\textbf{Übergang} & \textbf{Stelle} & \textbf{Übergang} & \textbf{Stelle}\\
\midrule
$000 \to 001$ & $3$ & $110 \to 111$ & $3$\\
$001 \to 011$ & $2$ & $111 \to 101$ & $2$\\
$011 \to 010$ & $3$ & $101 \to 100$ & $3$\\
$010 \to 110$ & $1$ & $100 \to 000$ & $1$\\
\bottomrule
\end{tabular}}
\end{center}
Acht Übergänge, jeder ändert genau eine Stelle. \checkmark

\begin{center}
\begin{minipage}{0.88\textwidth}
\kursiv{Gray-Codes werden bei Drehgebern verwendet --- etwa in der Winkelmessung. Beim Übergang von einer Stellung zur nächsten ändert sich nur ein einziges Bit; ein Ablesefehler kann daher höchstens eine Stellung danebenliegen. Bei gewöhnlicher Binärzählung geht $011 \to 100$ über drei Bits gleichzeitig, und ein ungünstiger Ablesezeitpunkt liefert einen völlig falschen Wert.}
\end{minipage}
\end{center}
\end{merksatz}

\section{Der Petersen-Graph ist nicht hamiltonsch}

Jetzt die versprochene Abrechnung mit dem Petersen-Graphen. Er erfüllt die Komponentenbedingung, hat keine Schnittecke, ist $3$-regulär und hochsymmetrisch --- und hat trotzdem keinen Hamiltonkreis. Der Beweis ist der erste in diesem Band, der eine Fallunterscheidung über mehrere Seiten braucht.

\begin{merksatz}
\boxtitel{Die Bezeichnungen}
\begin{center}
\begin{tikzpicture}[scale=1.5]
\foreach \i in {0,...,4} {
  \node[knotenb] (a\i) at (90+72*\i:1.35) {$a_{\i}$};
  \node[knotenb] (b\i) at (90+72*\i:0.68) {$b_{\i}$};
  \draw[kante] (a\i) -- (b\i);
}
\foreach \i/\j in {0/1,1/2,2/3,3/4,4/0} \draw[kante] (a\i) -- (a\j);
\foreach \i/\j in {0/2,1/3,2/4,3/0,4/1} \draw[kante] (b\i) -- (b\j);
\end{tikzpicture}
\end{center}

\begin{center}
\renewcommand{\arraystretch}{1.45}
\begin{tabular}{@{}P{2.4cm}P{3.0cm}P{2.2cm}@{}}
\toprule
\textbf{Sorte} & \textbf{Kanten} & \textbf{Anzahl}\\
\midrule
Außenkreis & $a_i a_{i+1}$ & $5$\\
Speichen & $a_i b_i$ & $5$\\
Innenstern & $b_i b_{i+2}$ & $5$\\
\bottomrule
\end{tabular}
\end{center}
Alle Indizes werden modulo $5$ gerechnet: $a_5$ bedeutet $a_0$.
\end{merksatz}

\begin{merksatz}
\boxtitel{Hilfssatz 7.7 --- Kreise überqueren einen Schnitt gerade oft}
Sei $C$ ein Kreis in einem Graphen und $A \subseteq V$. Dann ist die Anzahl der Kanten von $C$, die eine Ecke in $A$ und eine außerhalb haben, gerade.

\schritt{Beweis}
Wir laufen den Kreis einmal ab und notieren bei jedem Schritt, ob wir in $A$ oder außerhalb sind. Jede Überquerungskante wechselt die Seite, jede andere nicht.

Da wir am Ende wieder am Ausgangspunkt sind --- also auf derselben Seite wie zu Beginn ---, muss die Anzahl der Seitenwechsel gerade sein. $\blacksquare$
\end{merksatz}

\begin{merksatz}
\boxtitel{Satz 7.8 --- Der Petersen-Graph hat keinen Hamiltonkreis}
\schritt{Annahme}
Sei $H$ ein Hamiltonkreis. Er hat $10$ Ecken und daher $10$ Kanten, und jede Ecke hat in $H$ genau den Grad $2$.

\schritt{Schritt 1 --- Wie viele Speichen benutzt $H$?}
Sei $A = \{a_0, \ldots, a_4\}$ die Außenmenge. Die Kanten zwischen $A$ und dem Rest sind genau die fünf Speichen. Nach Hilfssatz 7.7 benutzt $H$ eine \kursiv{gerade} Anzahl davon. Sei $s$ diese Anzahl; also $s \in \{0, 2, 4\}$.

\schritt{Schritt 2 --- Die Kantenbilanz}
Wir zählen die Grade der Außenecken in $H$. Jede hat Grad $2$, also
\[ \sum_{i} \deg_H(a_i) = 10 . \]
Andererseits setzt sich diese Summe zusammen aus den Enden der benutzten Außenkanten (jede liefert $2$) und den Enden der benutzten Speichen (jede liefert $1$ an der Außenseite):
\[ 2 \cdot |{\text{Außenkanten in } H}| + s = 10 . \]
Also
\[ |{\text{Außenkanten in } H}| = \frac{10 - s}{2} . \]
Dieselbe Rechnung für die Innenecken ergibt denselben Wert für die Innenkanten.

\begin{center}
\renewcommand{\arraystretch}{1.45}
\begin{tabular}{@{}P{1.4cm}P{2.4cm}P{2.4cm}P{1.6cm}@{}}
\toprule
$s$ & \textbf{Außenkanten} & \textbf{Innenkanten} & \textbf{Summe}\\
\midrule
$0$ & $5$ & $5$ & $10$\\
$2$ & $4$ & $4$ & $10$\\
$4$ & $3$ & $3$ & $10$\\
\bottomrule
\end{tabular}
\end{center}
In allen Fällen kommt $10$ heraus, wie es sein muss.
\end{merksatz}

\begin{merksatz}
\boxtitel{Fall $s = 0$}
Dann benutzt $H$ alle fünf Außenkanten und alle fünf Innenkanten, aber keine Speiche.

Die fünf Außenkanten bilden den Kreis $a_0 a_1 a_2 a_3 a_4 a_0$ --- einen geschlossenen Kreis auf nur fünf Ecken. $H$ enthielte damit einen echten Teilkreis, was für einen Kreis mit zehn Ecken unmöglich ist: Ein Kreis hat an jeder Ecke genau zwei Kanten, und wenn fünf Ecken bereits untereinander geschlossen sind, kann keine Kante mehr nach außen führen.

Also ist $s = 0$ ausgeschlossen.
\end{merksatz}

\begin{merksatz}
\boxtitel{Fall $s = 2$}
$H$ benutzt $4$ der $5$ Außenkanten, $4$ der $5$ Innenkanten und $2$ Speichen.

\schritt{Was die Außenkanten bilden}
Vier der fünf Kanten des Kreises $a_0 \ldots a_4$: Das ist der Kreis ohne eine Kante, also ein \kursiv{Weg} durch alle fünf Außenecken. Sei $\{a_k, a_{k+1}\}$ die weggelassene Kante; dann sind die beiden Endecken des Weges genau $a_k$ und $a_{k+1}$.

\schritt{Wo die Speichen ansetzen müssen}
Die Ecken im Inneren des Weges haben schon Grad $2$ durch Außenkanten. Die beiden Endecken $a_k$ und $a_{k+1}$ haben erst Grad $1$ und brauchen je eine weitere Kante --- das können nur Speichen sein.

Die beiden benutzten Speichen sind also $\{a_k, b_k\}$ und $\{a_{k+1}, b_{k+1}\}$.

\schritt{Was innen passieren muss}
Nach demselben Argument bilden die vier Innenkanten einen Weg durch alle fünf Innenecken, und seine beiden Endecken müssen genau die von den Speichen getroffenen sein: $b_k$ und $b_{k+1}$.

\schritt{Der Widerspruch}
Der Innenstern ist der Kreis $b_0 b_2 b_4 b_1 b_3 b_0$. Lässt man daraus eine Kante weg, etwa $\{b_j, b_{j+2}\}$, so sind die Endecken des entstehenden Weges genau $b_j$ und $b_{j+2}$ --- zwei Ecken, deren Indizes sich um $2$ unterscheiden.

Gefordert waren aber $b_k$ und $b_{k+1}$ --- Indizes, die sich um $1$ unterscheiden. Es müsste also
\[ 1 \equiv \pm 2 \pmod 5 \]
gelten. Das ist falsch: $1 \neq 2$ und $1 \neq 3$ modulo $5$.

Also ist $s = 2$ ausgeschlossen.
\end{merksatz}

\begin{merksatz}
\boxtitel{Fall $s = 4$}
$H$ benutzt $3$ Außenkanten, $3$ Innenkanten und $4$ Speichen. Genau eine Speiche fehlt; sei es $\{a_k, b_k\}$. Wir setzen ohne Einschränkung $k = 0$ --- der Graph ist symmetrisch genug, dass man die Nummerierung so drehen darf.

\schritt{Schritt 1 --- Was an $a_0$ und $b_0$ geschieht}
$a_0$ benutzt keine Speiche, hat aber Grad $2$. Die einzigen anderen Kanten an $a_0$ sind $\{a_4, a_0\}$ und $\{a_0, a_1\}$. Beide müssen also in $H$ liegen.

Ebenso muss $b_0$ beide Innenkanten benutzen: $\{b_0, b_2\}$ und $\{b_0, b_3\}$.

\schritt{Schritt 2 --- Die übrigen Außenecken}
$a_1$ benutzt die Speiche $\{a_1,b_1\}$ und hat bereits $\{a_0,a_1\}$. Damit ist Grad $2$ erreicht --- die Kante $\{a_1, a_2\}$ liegt \kursiv{nicht} in $H$.

Ebenso hat $a_4$ die Speiche und $\{a_4,a_0\}$; also liegt $\{a_3,a_4\}$ nicht in $H$.

\schritt{Schritt 3 --- Die dritte Außenkante ist erzwungen}
Bisher zwei Außenkanten ($a_4a_0$ und $a_0a_1$); es müssen drei sein. Von den restlichen drei Außenkanten sind $a_1a_2$ und $a_3a_4$ ausgeschlossen. Bleibt nur $\{a_2, a_3\}$.

Kontrolle: $a_2$ hat dann Speiche und $\{a_2,a_3\}$ --- Grad $2$. \checkmark\ $a_3$ ebenso. \checkmark

\schritt{Schritt 4 --- Die übrigen Innenecken}
$b_2$ hat bereits $\{b_0,b_2\}$ und die Speiche $\{a_2,b_2\}$ --- Grad $2$ erreicht. Ebenso $b_3$.

Bleiben $b_1$ und $b_4$: Beide haben ihre Speiche und brauchen je genau eine Innenkante. Und es fehlt genau eine Innenkante an der Zahl $3$.

Diese eine Kante muss also $b_1$ mit $b_4$ verbinden. Ist $\{b_1, b_4\}$ überhaupt eine Kante? Der Innenstern verbindet $b_i$ mit $b_{i\pm2}$; und $4 \equiv 1 - 2 \pmod 5$. Ja, sie existiert.

\schritt{Schritt 5 --- Alles ist festgelegt --- und falsch}
Es gibt also genau eine Möglichkeit. Die zehn Kanten von $H$ wären:
\begin{center}
\renewcommand{\arraystretch}{1.42}
\begin{tabular}{@{}P{2.4cm}P{5.8cm}@{}}
\toprule
Außen & $a_4a_0$, $a_0a_1$, $a_2a_3$\\
Speichen & $a_1b_1$, $a_2b_2$, $a_3b_3$, $a_4b_4$\\
Innen & $b_0b_2$, $b_0b_3$, $b_1b_4$\\
\bottomrule
\end{tabular}
\end{center}

Jede Ecke hat darin Grad $2$ --- das lässt sich nachprüfen. Aber ist es ein \kursiv{Kreis}?

Wir laufen los bei $a_0$:
\[ a_0 \to a_1 \to b_1 \to b_4 \to a_4 \to a_0 . \]
Nach fünf Schritten sind wir zurück. Das ist ein Kreis der Länge $5$, nicht $10$.

Die übrigen fünf Ecken bilden den zweiten Kreis:
\[ a_2 \to a_3 \to b_3 \to b_0 \to b_2 \to a_2 . \]

\begin{center}
\begin{tikzpicture}[scale=1.5]
\foreach \i in {0,...,4} {
  \node[knotenb] (a\i) at (90+72*\i:1.35) {$a_{\i}$};
  \node[knotenb] (b\i) at (90+72*\i:0.68) {$b_{\i}$};
}
\foreach \i/\j in {0/1,1/2,2/3,3/4,4/0} \draw[kante, grau!35] (a\i) -- (a\j);
\foreach \i/\j in {0/2,1/3,2/4,3/0,4/1} \draw[kante, grau!35] (b\i) -- (b\j);
\foreach \i in {0,...,4} \draw[kante, grau!35] (a\i) -- (b\i);
\draw[kanted] (a4)--(a0); \draw[kanted] (a0)--(a1);
\draw[kanted] (a1)--(b1); \draw[kanted] (b1)--(b4); \draw[kanted] (b4)--(a4);
\draw[kante, gruen, line width=1.5pt] (a2)--(a3);
\draw[kante, gruen, line width=1.5pt] (a3)--(b3); \draw[kante, gruen, line width=1.5pt] (b3)--(b0);
\draw[kante, gruen, line width=1.5pt] (b0)--(b2); \draw[kante, gruen, line width=1.5pt] (b2)--(a2);
\end{tikzpicture}
\end{center}

Es entstehen zwei getrennte Fünfkreise, kein Hamiltonkreis. Also ist auch $s = 4$ ausgeschlossen.

\schritt{Schluss}
Alle drei möglichen Werte von $s$ führen auf einen Widerspruch. Der Petersen-Graph besitzt keinen Hamiltonkreis. $\blacksquare$
\end{merksatz}

\begin{merksatz}
\boxtitel{Und doch fast}
Der Petersen-Graph besitzt einen Hamilton\kursiv{weg}, zum Beispiel
\[ a_0 \to a_1 \to a_2 \to a_3 \to a_4 \to b_4 \to b_1 \to b_3 \to b_0 \to b_2 . \]
Prüfen Sie die Kanten: die vier Außenkanten, die Speiche $a_4b_4$, dann $b_4b_1$ (Differenz $2$), $b_1b_3$ (Differenz $2$), $b_3b_0$ (Differenz $2$), $b_0b_2$ (Differenz $2$). Alle vorhanden, alle zehn Ecken genau einmal. \checkmark

Mehr noch: Entfernt man aus dem Petersen-Graphen eine \kursiv{beliebige} Ecke, so ist der Rest hamiltonsch. Solche Graphen heißen \Begriff{hypohamiltonsch} --- sie sind es selbst nicht, aber jeder ihrer Eckenrestgraphen ist es. Der Petersen-Graph ist der kleinste.
\end{merksatz}

\section{Übungsaufgaben}

\begin{aufgabe}
\boxtitel{Aufgabe 7.1 --- Bipartite Graphen und Hamiltonkreise}
Zeigen Sie: Ist $G$ bipartit mit Bipartition $(A,B)$ und hamiltonsch, so gilt $|A| = |B|$. Folgern Sie, dass $K_{3,4}$ nicht hamiltonsch ist.

\loesung
\schritt{Schritt 1 --- Der Hamiltonkreis wechselt die Seiten}
Sei $H = (v_1, v_2, \ldots, v_n, v_1)$ ein Hamiltonkreis. Jede Kante von $H$ ist eine Kante von $G$, verbindet also eine Ecke aus $A$ mit einer aus $B$.

Aufeinanderfolgende Ecken liegen daher stets auf verschiedenen Seiten. Die Ecken wechseln entlang des Kreises also ab: $A, B, A, B, \ldots$

\schritt{Schritt 2 --- Der Kreis hat gerade Länge}
Nach Satz 2.12 enthält ein bipartiter Graph keinen ungeraden Kreis. Also ist $n$ gerade, sagen wir $n = 2r$.

\schritt{Schritt 3 --- Abzählen}
Bei striktem Abwechseln über $2r$ Ecken liegen genau $r$ Ecken in $A$ und $r$ in $B$. Da der Hamiltonkreis \kursiv{alle} Ecken enthält, ist
\[ |A| = r = |B| . \]

\schritt{Schritt 4 --- Anwendung auf $K_{3,4}$}
Hier ist $|A| = 3$ und $|B| = 4$, also $|A| \neq |B|$. Nach dem Bewiesenen ist $K_{3,4}$ nicht hamiltonsch.

\schritt{Gegenprobe über Satz 7.2}
Wählt man $S = A$ (die drei Ecken), so zerfällt $K_{3,4} - S$ in vier isolierte Ecken:
\[ c(G - S) = 4 \;>\; 3 = |S| . \]
Auch die Komponentenbedingung ist verletzt. Beide Wege führen zum selben Ergebnis. \checkmark

\schritt{Nachbemerkung}
$K_{4,4}$ dagegen \kursiv{ist} hamiltonsch: $a_1 b_1 a_2 b_2 a_3 b_3 a_4 b_4 a_1$. Die Bedingung $|A| = |B|$ ist bei vollständig bipartiten Graphen also nicht nur notwendig, sondern auch hinreichend.
\end{aufgabe}

\begin{aufgabe}
\boxtitel{Aufgabe 7.2 --- Ore prüfen}
Sei $G$ der Graph auf $\{1,\ldots,6\}$ mit den Kanten
\[ 12,\ 13,\ 14,\ 23,\ 25,\ 34,\ 36,\ 45,\ 46,\ 56 . \]
Prüfen Sie die Bedingungen von Dirac und Ore und entscheiden Sie über Hamiltonizität.

\loesung
\schritt{Schritt 1 --- Grade bestimmen}
\begin{center}
\renewcommand{\arraystretch}{1.42}
\begin{tabular}{@{}P{1.2cm}P{3.4cm}P{1.4cm}@{}}
\toprule
\textbf{Ecke} & \textbf{Nachbarn} & $\deg$\\
\midrule
$1$ & $2,3,4$ & $3$\\
$2$ & $1,3,5$ & $3$\\
$3$ & $1,2,4,6$ & $4$\\
$4$ & $1,3,5,6$ & $4$\\
$5$ & $2,4,6$ & $3$\\
$6$ & $3,4,5$ & $3$\\
\bottomrule
\end{tabular}
\end{center}
Probe: $3+3+4+4+3+3 = 20 = 2 \cdot 10$. \checkmark

\schritt{Schritt 2 --- Dirac}
Verlangt: $\deg(v) \ge n/2 = 3$ für alle $v$. Der kleinste Grad ist $3$. \textbf{Dirac ist erfüllt.}

\schritt{Schritt 3 --- Also ist $G$ hamiltonsch}
Nach Folgerung 7.4 besitzt $G$ einen Hamiltonkreis. Das Ergebnis steht damit schon fest; wir suchen ihn trotzdem.

\schritt{Schritt 4 --- Ore zur Kontrolle}
Die nicht benachbarten Paare sind $15$, $16$, $24$, $26$, $35$. Wir prüfen jedes:
\begin{center}
\renewcommand{\arraystretch}{1.4}
\begin{tabular}{@{}P{1.6cm}P{2.8cm}P{1.6cm}@{}}
\toprule
\textbf{Paar} & $\deg + \deg$ & $\ge 6$?\\
\midrule
$1,5$ & $3+3 = 6$ & \checkmark\\
$1,6$ & $3+3 = 6$ & \checkmark\\
$2,4$ & $3+4 = 7$ & \checkmark\\
$2,6$ & $3+3 = 6$ & \checkmark\\
$3,5$ & $4+3 = 7$ & \checkmark\\
\bottomrule
\end{tabular}
\end{center}
Ore ist ebenfalls erfüllt --- wie es sein muss, da Dirac schärfer ist.

\schritt{Schritt 5 --- Einen Hamiltonkreis angeben}
Versuch: $1 \to 2 \to 5 \to 6 \to 3 \to 4 \to 1$.

Kantenprüfung: $12$ \checkmark, $25$ \checkmark, $56$ \checkmark, $63$ \checkmark (in der Liste als $36$), $34$ \checkmark, $41$ \checkmark (als $14$).

Alle sechs Ecken genau einmal, sechs Kanten. \checkmark
\end{aufgabe}

\begin{aufgabe}
\boxtitel{Aufgabe 7.3 --- Ein Graph, den keine der Bedingungen erfasst}
Der Kreis $C_8$ ist hamiltonsch. Zeigen Sie, dass weder Dirac noch Ore das erkennen, und erklären Sie, was das über die Reichweite solcher Sätze aussagt.

\loesung
\schritt{Schritt 1 --- Dirac}
In $C_8$ hat jede Ecke Grad $2$. Verlangt wäre $n/2 = 4$. Also $2 < 4$: \textbf{Dirac greift nicht.}

\schritt{Schritt 2 --- Ore}
Nehmen wir zwei gegenüberliegende Ecken, etwa $v_1$ und $v_5$. Sie sind nicht benachbart, und
\[ \deg(v_1) + \deg(v_5) = 2 + 2 = 4 \;<\; 8 = n . \]
\textbf{Ore greift ebenfalls nicht.}

\schritt{Schritt 3 --- Trotzdem hamiltonsch}
$C_8$ \kursiv{ist} sein eigener Hamiltonkreis: $v_1 v_2 \ldots v_8 v_1$ enthält alle acht Ecken genau einmal.

\schritt{Schritt 4 --- Was daraus folgt}
\begin{center}
\begin{minipage}{0.9\textwidth}
Dirac und Ore sind \kursiv{hinreichende}, nicht notwendige Bedingungen. Sie sagen: \enquote{Wenn dicht genug, dann ja.} Sie sagen nichts über dünne Graphen --- und die interessanten Fälle sind meist dünn.

Zahlenmäßig: Dirac verlangt mindestens $n^2/4$ Kanten. $C_8$ hat $8$. Für $n = 8$ verlangt Dirac $16$. Der Abstand ist ein Faktor $2$, und er wächst mit $n$: Bei $n = 100$ verlangt Dirac $2500$ Kanten, $C_{100}$ hat $100$.
\end{minipage}
\end{center}

\begin{center}
\begin{minipage}{0.86\textwidth}
\kursiv{Man kennt keine hinreichende Bedingung, die \kursiv{alle} hamiltonschen Graphen erfasst und zugleich leicht prüfbar ist --- denn eine solche wäre ein vollständiges Kriterium. Kapitel~15 erklärt, warum das vermutlich unmöglich ist.}
\end{minipage}
\end{center}
\end{aufgabe}

\begin{geschichte}
\boxtitelg{Ein Spielzeug, das sich nicht verkaufte}
\Person{William Rowan Hamilton} (1805--1865) entwickelte 1856 eine algebraische Struktur, die er \kursiv{Icosian Calculus} nannte. Um sie anschaulich zu machen, entwarf er ein Spiel: einen hölzernen Dodekaeder mit zwanzig nummerierten Ecken, an denen Stifte steckten. Die Aufgabe war, mit einer Schnur einen Rundgang durch alle zwanzig Ecken zu legen.

Hamilton verkaufte die Rechte 1859 für fünfundzwanzig Pfund an einen Londoner Spielehändler. Das Spiel war ein Ladenhüter --- es galt als zu leicht. Für den Dodekaeder ist es das auch.

\begin{center}
\begin{minipage}{0.9\textwidth}
\kursiv{Die Frage nach solchen Rundgängen war schon vorher gestellt worden. \Person{Thomas Kirkman} untersuchte 1855 --- ein Jahr vor Hamilton --- systematisch, welche Polyeder einen Rundgang durch alle Ecken zulassen. Sein Aufsatz erschien in den Philosophical Transactions und ist die eigentliche erste Behandlung des Problems. Der Name blieb dennoch an Hamilton hängen.}
\end{minipage}
\end{center}

Der Fairness halber: Hamiltons Beitrag lag weniger im Problem als im algebraischen Rahmen, aus dem es entstand. Der Icosian Calculus ist ein früher Vorläufer der Gruppentheorie --- Band 11, Kapitel~3.
\end{geschichte}

\begin{forschung}
\boxtitelp{Warum es vermutlich kein Kriterium gibt}
Die Frage \enquote{Hat dieser Graph einen Hamiltonkreis?} gehört zu den ersten Problemen, für die \Person{Richard Karp} 1972 die \Begriff{NP-Vollständigkeit} nachwies. Was das bedeutet, entwickelt Band 3 der Informatikreihe; die kurze Fassung lautet:

\begin{center}
\renewcommand{\arraystretch}{1.5}
\begin{tabular}{@{}P{3.0cm}P{5.6cm}@{}}
\toprule
\textbf{Prüfen} & Ist ein Kreis vorgelegt, kann man in wenigen Schritten nachrechnen, ob er passt.\\
\textbf{Finden} & Kein Verfahren bekannt, das ihn in vertretbarer Zeit findet.\\
\textbf{Die Kopplung} & Gäbe es eines, so gäbe es zugleich schnelle Verfahren für Tausende anderer Probleme --- das ist die Frage \enquote{P gleich NP?}.\\
\bottomrule
\end{tabular}
\end{center}

Ein einfaches Kriterium vom Typ des Satzes 6.4 wäre genau so ein Verfahren. Wer eines fände, hätte damit eines der sieben Millenniumsprobleme gelöst.

\begin{center}
\begin{minipage}{0.88\textwidth}
\kursiv{Die Forschung geht deshalb andere Wege. Man untersucht eingeschränkte Graphenklassen, in denen es doch geht, und man untersucht \kursiv{typische} Graphen: In einem Zufallsgraphen mit $n$ Ecken ist ein Hamiltonkreis mit hoher Wahrscheinlichkeit vorhanden, sobald keine Ecke mehr den Grad $0$ oder $1$ hat --- ein Ergebnis von Komlós und Szemerédi (1983), auf das Kapitel~14 zurückkommt. Der schwierige Fall ist also nicht der Regelfall, sondern die Ausnahme.}
\end{minipage}
\end{center}
\end{forschung}

\begin{merksatz}
\boxtitel{Was Sie jetzt können}
\begin{center}
\renewcommand{\arraystretch}{1.5}
\begin{tabular}{@{}P{3.0cm}P{5.6cm}@{}}
\toprule
\textbf{Der Unterschied} & Euler betrifft Kanten und ist gelöst; Hamilton betrifft Ecken und ist es nicht\\
\textbf{Satz 7.2} & Komponentenbedingung --- das Werkzeug zum Widerlegen\\
\textbf{Satz 7.3} & Ore: $\deg(u) + \deg(v) \ge n$ genügt\\
\textbf{Folgerung 7.4} & Dirac als Sonderfall\\
\textbf{Satz 7.6} & $Q_d$ ist hamiltonsch --- Gray-Codes\\
\textbf{Satz 7.8} & Der Petersen-Graph ist es nicht\\
\bottomrule
\end{tabular}
\end{center}

\begin{center}
\begin{minipage}{0.88\textwidth}
\kursiv{Damit endet Teil~III. Ab Kapitel~8 geht es nicht mehr ums Durchlaufen, sondern ums Zeichnen: Wann lässt sich ein Graph ohne Überkreuzungen aufs Papier bringen?}
\end{minipage}
\end{center}
\end{merksatz}
